\usepackage{upgreek}
\usepackage{amsmath}
\usepackage{amssymb}

\usepackage{xcolor}
\usepackage{fvextra}
\usepackage{float}
\usepackage{caption}
\usepackage{etoolbox}
\usepackage{ifthen}

\newcommand{\dashedbox}[1]{
  \begin{tikzpicture}
    \node[draw, dashed, rounded corners=5pt, inner sep=10pt] {
      \begin{minipage}{0.8\textwidth}
        #1
      \end{minipage}
    };
  \end{tikzpicture}
}

% Configuración para bloques de código en Quarto
\DefineVerbatimEnvironment{Highlighting}{Verbatim}{%
    commandchars=\\\{\},
    breaksymbolleft={},
    breaklines=true,
    breakanywhere=true,
    breakautoindent=true,
    breaksymbolindentleft=0pt,
    breaksymbolsepleft=0pt,
    breaksymbolindentright=0pt,
    breaksymbolsepright=0pt,
    frame=single,
    framesep=2mm,
    rulecolor=\color{gray!30},
    baselinestretch=0.8,
    fontsize=\small
}

% Permitir que los bloques de código se dividan entre páginas
\makeatletter
\def\verbatim@processline{%
  \hbox to\columnwidth{\the\verbatim@line\strut\hss}%
  \ifx\verbatim@line\verbatim@empty\else%
    \ifnum\prevgraf>23%
      \penalty-200%
      \prevgraf0%
    \fi%
  \fi%
}
\makeatother

% Configuración adicional para listados largos
\emergencystretch=3em